\chapter{Development Story and Lessons Learned}

This chapter records how Smart Relay evolved and what decisions mattered. It captures the wisdom gained from building a production-ready AI-assisted proxy client with a modern GUI.

\section{Milestones}
\begin{itemize}
    \item \textbf{Scaffold}: set up Rust workspace, Axum control plane, AI client, and config store.
    \item \textbf{AI integration}: Ollama client with prompts that embed the live config snapshot.
    \item \textbf{Validation}: added HTTP error mapping and config validation to avoid unsafe applies.
    \item \textbf{GUI foundation}: Integrated \texttt{egui} for cross-platform native GUI with cyberpunk theme.
    \item \textbf{Subscription parsing}: Implemented Base64 decoding and multi-protocol parsing (VMess, VLESS, Shadowsocks, Trojan, SOCKS5).
    \item \textbf{System proxy integration}: Windows registry-based proxy configuration with Global/PAC mode support.
    \item \textbf{Endpoint management}: Connection testing, batch operations, and right-click context menus.
    \item \textbf{Internationalization}: Chinese character support with proper font configuration.
    \item \textbf{Comprehensive network diagnostics}: Multi-layer testing (DNS, TCP, routing, packet-level analysis) with detailed error categorization.
\end{itemize}

\section{Design Choices}

\subsection{Architecture}
\begin{itemize}
    \item Axum + Tokio for async control-plane tasks across platforms.
    \item Serde + TOML for human-friendly config with easy round trips.
    \item Adapter abstraction to decouple decision layer from core engines (xray, sing-box).
    \item Library crate structure: separate binaries for CLI and GUI, shared core logic.
    \item Immediate mode GUI (\texttt{egui}) for simplicity and cross-platform consistency.
\end{itemize}

\subsection{State Management}
\begin{itemize}
    \item \textbf{Config persistence}: TOML files in platform-specific config directories.
    \item \textbf{Async GUI}: Tokio runtime injected into \texttt{eframe} for async operations.
    \item \textbf{Channel-based communication}: \texttt{mpsc::channel} for passing async results to GUI thread.
    \item \textbf{One-time initialization}: Font configuration and theme application done once per frame.
\end{itemize}

\section{Critical Lessons Learned}

\subsection{Async Operations in Immediate Mode GUI}

The biggest challenge was integrating async operations (network requests, endpoint testing) with \texttt{egui}'s immediate mode update loop. \textbf{Key insight}: \texttt{eframe} runs in a synchronous context, but async operations require a Tokio runtime.

\textbf{Solution pattern}:
\begin{verbatim}
// In main.rs: Create runtime once
let rt = Arc::new(Runtime::new().expect("Failed to create Tokio runtime"));

// Pass to App
app.set_runtime(rt.clone());

// In App methods: Use rt.spawn() instead of tokio::spawn()
let (tx, rx) = mpsc::channel();
rt.spawn(async move {
    let result = async_operation().await;
    tx.send(result).unwrap();
});

// In update loop: Check for messages
if let Ok(msg) = rx.try_recv() {
    // Update UI state
}
\end{verbatim}

\textbf{Lessons}:
\begin{itemize}
    \item Never call \texttt{tokio::spawn} directly from GUI code—always use the injected runtime.
    \item Use channels (\texttt{mpsc}) to communicate between async tasks and GUI thread.
    \item Process channel messages in the \texttt{update()} method using \texttt{try\_recv()} to avoid blocking.
    \item Store receivers as \texttt{Option<mpsc::Receiver<T>>} in App state for processing across frames.
\end{itemize}

\subsection{Chinese Character Rendering}

Initial implementation showed Chinese characters as rectangles. \textbf{Root cause}: \texttt{egui}'s default fonts don't include CJK glyphs.

\textbf{Solutions tried}:
\begin{enumerate}
    \item \textbf{egui-chinese-font crate}: Automatic system font detection (works on most systems).
    \item \textbf{Manual font loading}: Fallback to Windows system fonts (\texttt{C:\textbackslash Windows\textbackslash Fonts\textbackslash msyh.ttc}).
    \item \textbf{Embedded fonts}: For maximum compatibility, embed Noto Sans CJK or similar.
\end{enumerate}

\textbf{Best practice}: Always test with actual Unicode content (subscription URLs with Chinese endpoint names) during development, not just ASCII.

\subsection{Windows System Proxy Integration}

Windows proxy configuration requires registry manipulation. \textbf{Key discovery}: Using \texttt{winreg} crate is more reliable than \texttt{netsh} commands.

\textbf{Implementation pattern}:
\begin{verbatim}
use winreg::enums::*;
use winreg::RegKey;

let hkcu = RegKey::predef(HKEY_CURRENT_USER);
let internet_settings = hkcu.open_subkey_with_flags(
    r"Software\Microsoft\Windows\CurrentVersion\Internet Settings",
    KEY_WRITE,
)?;

// Global proxy
internet_settings.set_value("ProxyEnable", &1u32)?;
internet_settings.set_value("ProxyServer", &proxy_str)?;
internet_settings.set_value("ProxyOverride", &exceptions)?;

// PAC mode
internet_settings.set_value("ProxyEnable", &0u32)?;
internet_settings.set_value("AutoConfigURL", &pac_url)?;
\end{verbatim}

\textbf{Lessons}:
\begin{itemize}
    \item Registry writes require \texttt{KEY\_WRITE} permission.
    \item Always provide user feedback when proxy changes are applied.
    \item Disable proxy on disconnect to restore user's original settings.
    \item Handle errors gracefully—proxy setup failures shouldn't crash the app.
\end{itemize}

\subsection{Subscription Parsing}

Subscription URLs return Base64-encoded content with multiple proxy protocols. \textbf{Challenges}:
\begin{itemize}
    \item Base64 decoding can fail—always handle errors.
    \item VMess JSON parsing requires handling multiple field names (\texttt{ps}, \texttt{name}, \texttt{remark}).
    \item Some protocols (VLESS, Trojan) have complex URL schemes.
    \item UTF-8 encoding must be preserved for Chinese endpoint names.
\end{itemize}

\textbf{Robust parsing pattern}:
\begin{verbatim}
// Try Base64 decode
let decoded = general_purpose::STANDARD.decode(base64_part)
    .map_err(|e| format!("Base64 decode failed: {}", e))?;

// Convert to UTF-8
let json_str = String::from_utf8(decoded)
    .map_err(|e| format!("UTF-8 decode failed: {}", e))?;

// Parse JSON with fallback field names
let name = json.get("ps")
    .or_else(|| json.get("name"))
    .or_else(|| json.get("remark"))
    .and_then(|v| v.as_str())
    .unwrap_or("Unknown Server");
\end{verbatim}

\subsection{Endpoint Testing and Connection Management}

Testing endpoints requires async TCP connections. \textbf{Pattern}:
\begin{verbatim}
pub async fn test_endpoint(host: &str, port: u16) -> Result<u64> {
    let start = Instant::now();
    match timeout(Duration::from_secs(5), 
                  TcpStream::connect((host, port))).await {
        Ok(Ok(_stream)) => Ok(start.elapsed().as_millis() as u64),
        Ok(Err(e)) => Err(anyhow!("Connection failed: {}", e)),
        Err(_) => Err(anyhow!("Connection timeout")),
    }
}
\end{verbatim}

\textbf{State management}:
\begin{itemize}
    \item Track testing state per endpoint (\texttt{testing: bool}).
    \item Store results in endpoint struct (\texttt{latency: Option<u64>}).
    \item Use visual indicators (spinner, color-coded latency).
    \item Batch testing: spawn multiple tests concurrently using \texttt{futures::join\_all}.
\end{itemize}

\subsection{Right-Click Context Menus}

\texttt{egui} doesn't have built-in context menus. \textbf{Solution}: Use \texttt{egui::Window} with fixed position.

\textbf{Pattern}:
\begin{verbatim}
// Detect right-click
if row_response.secondary_clicked() {
    self.right_click_menu = Some((
        row_response.interact_pointer_pos().unwrap(),
        idx
    ));
}

// Show menu window
if let Some((pos, idx)) = self.right_click_menu {
    egui::Window::new("context_menu")
        .fixed_pos(pos)
        .collapsible(false)
        .title_bar(false)
        .show(ui.ctx(), |ui| {
            // Menu items
        });
}
\end{verbatim}

\section{What Worked Well}
\begin{itemize}
    \item Small prompts with explicit sections lead to clearer AI responses.
    \item Early validation hooks prevented bad configs from reaching adapters.
    \item Keeping the control plane stateless (config on disk) simplifies restart/upgrade.
    \item \textbf{Channel-based async communication}: Clean separation between async tasks and GUI updates.
    \item \textbf{Default trait implementations}: Using \texttt{Default} for struct initialization reduces boilerplate.
    \item \textbf{Error handling}: Using \texttt{anyhow::Result} for error propagation simplifies debugging.
    \item \textbf{Visual feedback}: Color-coded latency, connection indicators, and status messages improve UX.
\end{itemize}

\section{Common Pitfalls and Solutions}

\subsection{Pitfall: "No reactor running" Error}
\textbf{Symptom}: \texttt{thread 'main' panicked: there is no reactor running}.\\
\textbf{Cause}: Calling \texttt{tokio::spawn} from non-async context.\\
\textbf{Solution}: Create \texttt{Runtime} in \texttt{main.rs} and pass to GUI via \texttt{Arc<Runtime>}.

\subsection{Pitfall: Moved Value Errors}
\textbf{Symptom}: \texttt{error[E0382]: use of moved value}.\\
\textbf{Cause}: Using a value in multiple async closures.\\
\textbf{Solution}: Clone values before moving into closures:
\begin{verbatim}
let host_clone = host.clone();
rt.spawn(async move {
    test_endpoint(&host_clone, port).await;
});
\end{verbatim}

\subsection{Pitfall: Font Configuration Not Applied}
\textbf{Symptom}: Chinese characters show as rectangles.\\
\textbf{Cause}: Fonts configured every frame, or default fonts don't support CJK.\\
\textbf{Solution}: Configure fonts once in \texttt{main.rs} using \texttt{egui-chinese-font}, or load system fonts manually.

\subsection{Pitfall: Grid Column Mismatch}
\textbf{Symptom}: UI layout breaks when adding new columns.\\
\textbf{Cause}: \texttt{num\_columns} doesn't match actual columns rendered.\\
\textbf{Solution}: Always update \texttt{num\_columns} when modifying grid structure.

\section{What to Improve Next}
\begin{itemize}
    \item Add schema-validated JSON outputs from AI to reduce parsing ambiguity.
    \item Implement real adapters and probe-driven scoring loops.
    \item Build integration tests for API endpoints and adapter lifecycles.
    \item Add CI to run \texttt{cargo fmt/clippy/test} and LaTeX build checks.
    \item \textbf{Test result persistence}: Store test results in config for faster UI updates.
    \item \textbf{Connection pooling}: Reuse TCP connections for batch testing.
    \item \textbf{Clipboard integration}: Copy endpoint URLs to clipboard from context menu.
    \item \textbf{Export/Import}: Save endpoint lists as JSON for backup/sharing.
\end{itemize}

\section{Performance Considerations}

\begin{itemize}
    \item \textbf{Subscription parsing}: Large subscriptions (80+ endpoints) parse in \textasciitilde{}200ms—acceptable for one-time import.
    \item \textbf{Endpoint testing}: 5-second timeout per endpoint; batch testing runs concurrently.
    \item \textbf{GUI responsiveness}: All async operations are non-blocking; UI remains interactive.
    \item \textbf{Memory}: Endpoint list stored in memory; consider pagination for 1000+ endpoints.
\end{itemize}

\subsection{Network Diagnostics and Deep Packet Analysis}

One of the most challenging aspects was implementing comprehensive endpoint testing that provides actionable diagnostics. The initial implementation only tested basic TCP connectivity, but users needed to understand \textit{why} connections failed.

\textbf{The Challenge:}
\begin{itemize}
    \item Simple TCP tests often fail even when endpoints work through proxies.
    \item Users need to understand if failures are due to DNS, routing, firewalls, or protocol issues.
    \item Different error types require different troubleshooting approaches.
\end{itemize}

\textbf{The Solution - Multi-Layer Testing:}
\begin{enumerate}
    \item \textbf{Layer 7 (DNS)}: Test hostname resolution first - if DNS fails, nothing else matters.
    \item \textbf{Layer 3 (Routing)}: Check routing table and traceroute to understand network path.
    \item \textbf{Layer 4 (TCP)}: Test TCP handshake with detailed packet flow analysis.
    \item \textbf{Layer 2/3 (Packets)}: Provide packet-level details (TCP headers, flags, sequence numbers).
    \item \textbf{Smart Testing}: Try multiple approaches (direct, via proxy, extended timeout).
\end{enumerate}

\textbf{Key Insights:}
\begin{itemize}
    \item \textbf{Connection Refused vs Timeout}: Refused means port is closed (RST received); timeout means no response (firewall dropping silently).
    \item \textbf{Packet Flow Visualization}: Showing "SYN → (no response)" helps users understand firewall behavior.
    \item \textbf{Error Categorization}: DNS failures, connection refused, timeouts, and network unreachable each require different fixes.
    \item \textbf{Diagnostic Depth}: The more detailed the diagnostics, the easier it is for users (and AI) to diagnose issues.
\end{itemize}

\textbf{Implementation Details:}
\begin{itemize}
    \item Used \texttt{tokio::net::lookup\_host} for DNS resolution with timeout.
    \item Used \texttt{tokio::net::TcpStream::connect} for TCP testing with configurable timeouts.
    \item Integrated system commands (\texttt{ipconfig}, \texttt{route}, \texttt{tracert}) for routing information.
    \item Created detailed diagnostic messages explaining packet flow and connection states.
    \item Categorized errors to provide actionable feedback (e.g., "Firewall dropping silently" vs "Port closed").
\end{itemize}

This multi-layer approach provides the context needed to understand VPN connection issues, which is especially important when endpoints may work through proxies but fail direct connections.

\section{Personal Checklist for Contributors}
\begin{itemize}
    \item Run \texttt{cargo test} before pushing changes.
    \item Keep prompts and configs under version control to aid rollbacks.
    \item Document decisions in \texttt{docs/dev\_story.md} to keep the narrative fresh.
    \item \textbf{Test with real data}: Use actual subscription URLs, not mock data.
    \item \textbf{Test on target platform}: Windows proxy features only work on Windows.
    \item \textbf{Handle errors gracefully}: Never panic on user input or network errors.
    \item \textbf{Provide feedback}: Show loading states, success messages, and error details.
    \item \textbf{Understand network layers}: When debugging connection issues, think about which OSI layer is failing.
    \item \textbf{Provide detailed diagnostics}: Users need context to understand why connections fail.
\end{itemize}

